
\subsection*{\SGAInlineRefTarget{sga3:i:subsection:4-4}4.4. Diagonalizable groups}

\medskip
\noindent\SGAInlineRefTarget{sga3:i:paragraph:4-4-1}\textbf{4.4.1.}
The construction of \(\mathbb G_m\) generalizes as follows. Let \(M\)
be a commutative group, and let \(M_{\mathbb Z}\) be the group scheme
associated with it by the procedure of \SGARefLink{sga3:i:subsection:4-1}{4.1}. Consider the functor defined
by
\[
  \mathbf D(M)(S)
    =\operatorname{Hom}_{\mathrm{groups}}(M,\mathbb G_m(S))
    \simeq
      \underline{\operatorname{Hom}}_{S\text{-}\mathbf{Gr}}
        (M_S,\mathbb G_{m,S}).
\]

This is a commutative \(\widehat{(\mathbf{Sch})}\)-group represented by a
group scheme denoted \(D(M)\); thus, by definition,
\[
  D(M)\simeq
  \underline{\operatorname{Hom}}_{(\mathbf{Sch})\text{-}\mathbf{Gr}}
    (M_{\mathbb Z},\mathbb G_m).
\]
Indeed, set
\[
  D(M)=\operatorname{Spec}\mathbb Z[M],
\]
where \(\mathbb Z[M]\) is the group algebra of \(M\) over \(\mathbb Z\).
Then
\[
  D(M)(S)
    =\operatorname{Hom}_{\mathrm{Alg.}}
       \bigl(\mathbb Z[M],\Gamma(S,\mathcal O_S)\bigr)
    \simeq
      \operatorname{Hom}_{\mathrm{groups}}
       \bigl(M,\Gamma(S,\mathcal O_S)^{\times}\bigr)
\]
by the very definition of the algebra \(\mathbb Z[M]\).

\medskip
\noindent\SGAInlineRefTarget{sga3:i:paragraph:4-4-2}\textbf{4.4.2.}
For every scheme \(S\), one therefore has an \(S\)-group scheme affine
over \(S\),
\[
  D_S(M)=D(M)_S
    =\bigl(
      \underline{\operatorname{Hom}}_{(\mathbf{Sch})\text{-}\mathbf{Gr}}
        (M_{\mathbb Z},\mathbb G_m)
      \bigr)_S
    =\underline{\operatorname{Hom}}_{S\text{-}\mathbf{Gr}}
       (M_S,\mathbb G_{m,S}).
\]
It is associated with the \(\mathcal O_S\)-bialgebra
\(\mathcal O_S[M]\), endowed with the diagonal map and augmentation
defined by
\[
  \Delta(x)=x\otimes x,
  \qquad
  \varepsilon(x)=1,
  \qquad x\in M.
\]

\medskip
\noindent\SGAInlineRefTarget{sga3:i:paragraph:4-4-3}\textbf{4.4.3.}
If \(f:M\to N\) is a homomorphism of commutative groups, one defines in
the evident way a morphism of \(S\)-groups
\[
  D_S(f):D_S(N)\longrightarrow D_S(M),
\]
and hence a functor
\[
  D_S:M\longmapsto D_S(M)
\]
from the category of abelian groups to the category of \(S\)-groups
affine over \(S\). It may also be defined as the composite of the
functor \(M\mapsto M_S\) and the functor
\[
  M_S\longmapsto
  \underline{\operatorname{Hom}}_{S\text{-}\mathbf{Gr}}
    (M_S,\mathbb G_{m,S}).
\]
This functor commutes with base extension.

An \(S\)-group isomorphic to a group of the form \(D_S(M)\) is said to
be \emph{diagonalizable}. Notice that the elements of \(M\) are
interpreted as certain characters of \(D_S(M)\), that is, as certain
elements of
\[
  \operatorname{Hom}_{S\text{-}\mathbf{Gr}}
    (D_S(M),\mathbb G_{m,S}).
\]
(Compare VIII, 1.)

\medskip
\noindent\SGAInlineRefTarget{sga3:i:paragraph:4-4-4}\textbf{4.4.4.}
Let us give some examples of diagonalizable groups. First,
\[
  D(\mathbb Z)=\mathbb G_m,
  \qquad
  D(\mathbb Z^r)=(\mathbb G_m)^r.
\]
Set
\[
  \boldsymbol\mu_n=D(\mathbb Z/n\mathbb Z),
\]
and call it the \emph{group of \(n\)-th roots of unity}. Indeed,
\[
  \boldsymbol\mu_n(S)
    =\operatorname{Hom}_{\mathrm{groups}}
       \bigl(\mathbb Z/n\mathbb Z,\Gamma(S,\mathcal O_S)^{\times}\bigr)
    =\{f\in\Gamma(S,\mathcal O_S)\mid f^n=1\}.
\]
The \(S\)-group \(\boldsymbol\mu_{n,S}\) corresponds to the
\(\mathcal O_S\)-algebra \(\mathcal O_S[T]/(T^n-1)\). Suppose, in
particular, that \(S\) is the spectrum of a field \(k\) of
characteristic \(p=n\). On setting \(T-1=s\), one obtains
\[
  k[T]/(T^p-1)=k[s]/(s^p),
\]
which shows that the underlying topological space of
\(\boldsymbol\mu_{p,k}\) consists of one point, whose local ring is the
Artinian \(k\)-algebra \(k[s]/(s^p)\). In the same vein, note that the
\(S\)-schemes \(\mathbb G_{a,S}\), \(\mathbb G_{m,S}\), and
\(\mathbf O_S\) are smooth over \(S\), and that \(D_S(M)\) is

flat over \(S\). Moreover, it is formally smooth (respectively, smooth)
over \(S\) if and only if no residue characteristic of \(S\) divides
the torsion of \(M\) (respectively, and if in addition \(M\) is of
finite type); see \SGARefLink{sga3:viii:proposition:2:1}{Expos\'e VIII, 2.1}.

\subsection*{\SGAInlineRefTarget{sga3:i:subsection:4-5}4.5. Other examples of groups}

The preceding procedure applies to the ``classical groups'' (linear
groups \(\operatorname{GL}_n\), symplectic groups
\(\operatorname{Sp}_n\), orthogonal groups \(\operatorname O_n\), and
so forth). For example, one defines \(\operatorname{GL}_n\) as
representing the functor \(\mathbf{GL}_n\) such that
\[
  \mathbf{GL}_n(S)
    =\operatorname{GL}\bigl(n,\Gamma(S,\mathcal O_S)\bigr)
    =\operatorname{Aut}_{\mathcal O_S}(\mathcal O_S^n).
\]
It may be constructed, for example, as the open subscheme of
\(\operatorname{Spec}\mathbb Z[T_{ij}]\), \(1\leq i,j\leq n\), defined
by the function
\[
  f=\det\bigl((T_{ij})_{i,j=1}^n\bigr),
\]
that is, as \(\operatorname{Spec}\mathbb Z[T_{ij},f^{-1}]\).

\subsection*{\SGAInlineRefTarget{sga3:i:subsection:4-6}4.6. Functor-modules in the category of schemes}

We shall associate with every \(\mathcal O_S\)-module on a scheme \(S\)
a \(\mathbf O_S\)-module, where \(\mathbf O_S\) denotes the functor-ring
introduced in 4.3.3. This can be done in two different ways, as follows.

\medskip
\noindent\SGAInlineRefTarget{sga3:i:definition:4-6-1}\textbf{Definition 4.6.1.}
Let \(S\) be a scheme. For every \(\mathcal O_S\)-module \(\mathscr F\),
denote by \(\mathbf V(\mathscr F)\) and \(\mathbf W(\mathscr F)\) the
contravariant functors on \((\mathbf{Sch})_{/S}\) defined by
\begin{align*}
  \mathbf V(\mathscr F)(S')
    &=\operatorname{Hom}_{\mathcal O_{S'}}
       (\mathscr F\otimes_{\mathcal O_S}\mathcal O_{S'},
        \mathcal O_{S'}),\\
  \mathbf W(\mathscr F)(S')
    &=\Gamma(S',\mathscr F\otimes_{\mathcal O_S}\mathcal O_{S'}),
\end{align*}
where \(\mathscr F\otimes_{\mathcal O_S}\mathcal O_{S'}\) denotes the
inverse image on \(S'\) of the \(\mathcal O_S\)-module \(\mathscr F\).

Then \(\mathbf V(\mathscr F)\) and \(\mathbf W(\mathscr F)\) are
endowed in the evident way with \(\mathbf O_S\)-module structures
(recall that
\(\mathbf O_S(S')=\Gamma(S',\mathcal O_{S'})
=\mathbf W(\mathcal O_S)(S')\)). Thus one has in fact defined two
functors \(\mathbf V\) and \(\mathbf W\) from the category of
\(\mathcal O_S\)-modules to the category of \(\mathbf O_S\)-modules,
with \(\mathbf V\) contravariant and \(\mathbf W\) covariant.

For the remainder of this paragraph, we restrict to the case in which
the \(\mathcal O_S\)-modules in question are quasi-coherent. Thus we
regard \(\mathbf V\) and \(\mathbf W\) as functors
\[
  \mathbf V:(\mathcal O_S\text{-}\mathbf{Mod.q.c.})^{\circ}
    \longrightarrow(\mathbf O_S\text{-}\mathbf{Mod.}),
  \qquad
  \mathbf W:(\mathcal O_S\text{-}\mathbf{Mod.q.c.})
    \longrightarrow(\mathbf O_S\text{-}\mathbf{Mod.}).
\]

\medskip
\noindent\SGAInlineRefTarget{sga3:i:remark:4-6-1-1}\textbf{Remark 4.6.1.1.}
The reader will notice that, in what follows, every proposition involving
only the functor \(\mathbf W\) remains valid for arbitrary
\(\mathcal O_S\)-modules, not necessarily quasi-coherent.

\medskip
\noindent\SGAInlineRefTarget{sga3:i:proposition:4-6-2}\textbf{Proposition 4.6.2.}
\begin{enumerate}
  \item[\SGAInlineRefTarget{sga3:i:proposition:4-6-2:item:i}(i)]
  The functors \(\mathbf V\) and \(\mathbf W\) commute with base
  extension: if \(S'\) lies over \(S\) and \(\mathscr F\) is a
  quasi-coherent \(\mathcal O_S\)-module, then
  \[
    \mathbf V(\mathscr F\otimes\mathcal O_{S'})
      \simeq\mathbf V(\mathscr F)_{S'},
    \qquad
    \mathbf W(\mathscr F\otimes\mathcal O_{S'})
      \simeq\mathbf W(\mathscr F)_{S'}.
  \]

  \item[\SGAInlineRefTarget{sga3:i:proposition:4-6-2:item:ii}(ii)]
  The functors \(\mathbf V\) and \(\mathbf W\) are fully faithful: the
  canonical maps
  \[
    \operatorname{Hom}_{\mathcal O_S}(\mathscr F,\mathscr F')
      \longrightarrow
    \operatorname{Hom}_{\mathbf O_S}
      (\mathbf V(\mathscr F'),\mathbf V(\mathscr F)),
  \]
  \[
    \operatorname{Hom}_{\mathcal O_S}(\mathscr F,\mathscr F')
      \longrightarrow
    \operatorname{Hom}_{\mathbf O_S}
      (\mathbf W(\mathscr F),\mathbf W(\mathscr F'))
  \]
  are bijective.

  \item[\SGAInlineRefTarget{sga3:i:proposition:4-6-2:item:iii}(iii)]
  The functors \(\mathbf V\) and \(\mathbf W\) are additive:
  \[
    \mathbf V(\mathscr F\oplus\mathscr F')
      \simeq
    \mathbf V(\mathscr F)\times_S\mathbf V(\mathscr F'),
    \qquad
    \mathbf W(\mathscr F\oplus\mathscr F')
      \simeq
    \mathbf W(\mathscr F)\times_S\mathbf W(\mathscr F').
  \]
\end{enumerate}

Parts \SGARefLink{sga3:i:proposition:4-6-2:item:i}{(i)} and \SGARefLink{sga3:i:proposition:4-6-2:item:iii}{(iii)} are immediate from the definitions. For \SGARefLink{sga3:i:proposition:4-6-2:item:ii}{(ii)}, take
the \(S'\) to be open subschemes of \(S\). We leave the proof to the
reader (for \(\mathbf V\), use EGA~II, 1.7.14).

Recall (cf. \SGARefLink{sga3:i:definition:3-1-4}{3.1.4}) that if \(\mathbf F\) and \(\mathbf F'\) are
\(\mathbf O_S\)-modules, then
\(\underline{\operatorname{Hom}}_{\mathbf O_S}(\mathbf F,\mathbf F')\)
denotes the \(S\)-functor in abelian groups that assigns to every
\(S'\to S\) the group
\[
  \operatorname{Hom}_{\mathbf O_{S'}}(\mathbf F_{S'},\mathbf F'_{S'}).
\]

\medskip
\noindent\SGAInlineRefTarget{sga3:i:proposition:4-6-3}\textbf{Proposition 4.6.3.}
\textit{There are canonical morphisms in
\((\mathbf O_S\text{-}\mathbf{Mod.})\):}
\footnote{Editor's note: the isomorphism
\(\underline{\operatorname{Hom}}_{\mathbf O_S}
(\mathbf W(\mathscr F),\mathbf W(\mathscr F'))
\simeq
\underline{\operatorname{Hom}}_{\mathbf O_S}
(\mathbf V(\mathscr F'),\mathbf V(\mathscr F))\)
has been added.}

\begin{figure}[H]
  \centering
  \begin{tikzcd}[column sep=3.2em,row sep=3.5em]
  \underline{\operatorname{Hom}}_{\mathbf O_S}
    \bigl(\mathbf W(\mathscr F),\mathbf W(\mathscr F')\bigr)
    \arrow[rr,equal]
  &&
  \underline{\operatorname{Hom}}_{\mathbf O_S}
    \bigl(\mathbf V(\mathscr F'),\mathbf V(\mathscr F)\bigr)
  \\
  &
  \mathbf W\!\left(
    \mathscr Hom_{\mathcal O_S}(\mathscr F,\mathscr F')
  \right)
    \arrow[ul]
    \arrow[ur]
  &
  {} .
\end{tikzcd}

  \phantomsection\label{fig:sga3-exp1-p026-prop463-canonical-morphisms}
\end{figure}

This follows immediately from \SGARefLink{sga3:i:proposition:4-6-2:item:i}{4.6.2(i)} and \SGARefLink{sga3:i:proposition:4-6-2:item:ii}{(ii)}.

\medskip
\noindent\SGAInlineRefTarget{sga3:i:notation:4-6-3-1}\textbf{Notation 4.6.3.1.}
Let \(\mathscr F\) be a quasi-coherent \(\mathcal O_S\)-module. By
EGA~II, 1.7.8, the \(S\)-functor \(\mathbf V(\mathscr F)\) is
represented by an \(S\)-scheme affine over \(S\), denoted
\(\mathbb V(\mathscr F)\), and called the \emph{vector fibration}
defined by \(\mathscr F\):
\[
  \mathbb V(\mathscr F)=\operatorname{Spec}(\mathcal S(\mathscr F)),
\]
where \(\mathcal S(\mathscr F)\) denotes the symmetric algebra of the
\(\mathcal O_S\)-module \(\mathscr F\).

\medskip
\noindent\SGAInlineRefTarget{sga3:i:proposition:4-6-4}\textbf{Proposition 4.6.4.}
\textit{Let \(\mathscr F\) and \(\mathscr F'\) be two quasi-coherent
\(\mathcal O_S\)-modules, and let \(\mathscr A\) be a quasi-coherent
\(\mathcal O_S\)-algebra. There is a functorial isomorphism
\[
  \operatorname{Hom}_S
   \bigl(\operatorname{Spec}(\mathscr A),
     \underline{\operatorname{Hom}}_{\mathbf O_S}
       (\mathbf W(\mathscr F'),\mathbf W(\mathscr F))\bigr)
  \xrightarrow{\ \sim\ }
  \operatorname{Hom}_{\mathcal O_S}
    (\mathscr F',\mathscr F\otimes_{\mathcal O_S}\mathscr A).
\]}

Indeed, put \(X=\operatorname{Spec}(\mathscr A)\). The left-hand side is
canonically isomorphic to
\[
  \underline{\operatorname{Hom}}_{\mathbf O_S}
    (\mathbf W(\mathscr F'),\mathbf W(\mathscr F))(X),
\]
that is, by definition, to
\begin{align*}
  &\operatorname{Hom}_{\mathbf O_X}
    (\mathbf W(\mathscr F')_X,\mathbf W(\mathscr F)_X)\\
  &\qquad\simeq
  \operatorname{Hom}_{\mathbf O_X}
    \bigl(\mathbf W(\mathscr F'\otimes\mathcal O_X),
          \mathbf W(\mathscr F\otimes\mathcal O_X)\bigr)
\end{align*}
(cf. \SGARefLink{sga3:i:proposition:4-6-2:item:i}{4.6.2(i)}). By \SGARefLink{sga3:i:proposition:4-6-2:item:ii}{4.6.2(ii)}, this may also be written as
\[
  \operatorname{Hom}_{\mathcal O_X}
    (\mathscr F'\otimes\mathcal O_X,\mathscr F\otimes\mathcal O_X)
  =
  \operatorname{Hom}_{\mathcal O_S}
    (\mathscr F',\pi_*\pi^*(\mathscr F)),
\]
where \(\pi:X\to S\) is the structural morphism. But EGA~II, 1.4.7
gives
\[
  \pi_*\pi^*(\mathscr F)\simeq\mathscr F\otimes\mathscr A,
\]
which completes the proof.

\medskip
\noindent\SGAInlineRefTarget{sga3:i:corollary:4-6-4-1}\textbf{Corollary 4.6.4.1.}
\textit{There is a canonical isomorphism
\[
  \mathbf W(\mathscr F\otimes\mathscr A)
    \simeq
  \underline{\operatorname{Hom}}_S
    (\operatorname{Spec}(\mathscr A),\mathbf W(\mathscr F)).
\]}

Indeed,
\footnote{Editor's note: the original has been expanded in what follows.}
let \(f:S'\to S\) be an \(S\)-scheme and put
\(X'=X\times_S S'\). There is a Cartesian square

\begin{figure}[H]
  \centering
  \begin{tikzcd}[column sep=4.0em,row sep=3.0em]
  X'
    \arrow[r,"f'"]
    \arrow[d,"\pi'"']
  &
  X
    \arrow[d,"\pi"]
  \\
  S'
    \arrow[r,"f"]
  &
  S
\end{tikzcd}

  \phantomsection\label{fig:sga3-exp1-p027-cor4641-cartesian-square}
\end{figure}

By EGA~II, 1.5.2, \(X'\) is affine over \(S'\) and
\(\pi'_*\mathcal O_{X'}=f^*(\mathscr A)\). Hence
\begin{align*}
  &\underline{\operatorname{Hom}}_S
     (\operatorname{Spec}(\mathscr A),\mathbf W(\mathscr F))(S')\\
  &\qquad=
  \operatorname{Hom}_{S'}
     \bigl(\operatorname{Spec}(f^*\mathscr A),
           \mathbf W(f^*\mathscr F)\bigr).
\end{align*}
By \SGARefLink{sga3:i:proposition:4-6-4}{4.6.4}, applied to \(f^*\mathscr F\),
\(\mathscr F'=\mathcal O_{S'}\), and \(f^*\mathscr A\), this is equal
to
\begin{align*}
  \Gamma(S',f^*\mathscr F\otimes f^*\mathscr A)
    &=\Gamma(S',f^*(\mathscr F\otimes\mathscr A))\\
    &=\mathbf W(\mathscr F\otimes\mathscr A)(S').
\end{align*}

\medskip
\noindent\SGAInlineRefTarget{sga3:i:proposition:4-6-5}\textbf{Proposition 4.6.5.}
\textit{If \(\mathscr F\) and \(\mathscr F'\) are two
\(\mathcal O_S\)-modules locally free of finite type, the morphisms of
\SGARefLink{sga3:i:proposition:4-6-3}{4.6.3} are isomorphisms.}

Indeed, for every \(S'\to S\),
\begin{align*}
  &\mathbf W
    \bigl(\mathscr{H}om_{\mathcal O_S}(\mathscr F,\mathscr F')\bigr)(S')\\
  &\quad=
    \Gamma\bigl(S',
      \mathscr{H}om_{\mathcal O_S}(\mathscr F,\mathscr F')
        \otimes\mathcal O_{S'}\bigr)\\
  &\quad=
    \operatorname{Hom}_{\mathcal O_{S'}}
      (\mathscr F\otimes\mathcal O_{S'},
       \mathscr F'\otimes\mathcal O_{S'}).
\end{align*}
The last member is isomorphic both to
\[
  \underline{\operatorname{Hom}}_{\mathbf O_S}
    (\mathbf W(\mathscr F),\mathbf W(\mathscr F'))(S')
\]
and to
\[
  \underline{\operatorname{Hom}}_{\mathbf O_S}
    (\mathbf V(\mathscr F'),\mathbf V(\mathscr F))(S')
\]
by \SGARefLink{sga3:i:proposition:4-6-2:item:i}{4.6.2(i)} and \SGARefLink{sga3:i:proposition:4-6-2:item:ii}{(ii)}.

\medskip
\noindent\SGAInlineRefTarget{sga3:i:corollary:4-6-5-1}\textbf{Corollary 4.6.5.1.}
\textit{Let \(\mathscr F\) be an \(\mathcal O_S\)-module locally free
of finite type. Put
\(\mathscr F^{\vee}
=\mathscr{H}om_{\mathcal O_S}(\mathscr F,\mathcal O_S)\). There are
canonical isomorphisms
\[
  \mathbf W(\mathscr F^{\vee})
    \simeq
  \underline{\operatorname{Hom}}_{\mathbf O_S}
    (\mathbf W(\mathscr F),\mathbf O_S)
    \simeq\mathbf V(\mathscr F),
\]
\[
  \mathbf V(\mathscr F^{\vee})
    \simeq
  \underline{\operatorname{Hom}}_{\mathbf O_S}
    (\mathbf V(\mathscr F),\mathbf O_S)
    \simeq\mathbf W(\mathscr F).
\]}

Finally, we have the following proposition.

\medskip
\noindent\SGAInlineRefTarget{sga3:i:proposition:4-6-6}\textbf{Proposition 4.6.6.}
\textit{Let \(f:\mathscr F\to\mathscr F'\) be a morphism of
\(\mathcal O_S\)-modules locally free of finite rank. The morphism
\[
  \mathbf W(f):\mathbf W(\mathscr F)\longrightarrow\mathbf W(\mathscr F')
\]
is a monomorphism if and only if \(f\) identifies \(\mathscr F\)
locally with a direct summand of \(\mathscr F'\).}

The direct assertion is essentially contained in EGA~0\textsubscript{I},
5.5.5.
\footnote{Editor's note: the following sentence has been expanded.}
Conversely, if \(\mathscr F\) is locally a direct summand of
\(\mathscr F'\), then for every \(\pi:S'\to S\), the module
\(\pi^*\mathscr F\) is a submodule of \(\pi^*\mathscr F'\), since this
holds locally as a direct summand. Thus
\[
  \mathbf W(\mathscr F)(S')=\Gamma(S',\pi^*\mathscr F)
\]
is a submodule of
\[
  \mathbf W(\mathscr F')(S')=\Gamma(S',\pi^*\mathscr F').
\]

\subsection*{\SGAInlineRefTarget{sga3:i:subsection:4-7}4.7. The category of \(G\)-\(\mathcal O_S\)-modules}

Let \(G\) be an \(S\)-group and \(\mathscr F\) an
\(\mathcal O_S\)-module; \(\mathbf W(\mathscr F)\) has a
\(\mathbf O_S\)-module structure.

\medskip
\noindent\SGAInlineRefTarget{sga3:i:definition:4-7-1}\textbf{Definition 4.7.1.}
A \emph{\(G\)-\(\mathcal O_S\)-module structure} on \(\mathscr F\) is
an \(\mathbf h_G\)-\(\mathbf O_S\)-module structure on
\(\mathbf W(\mathscr F)\) (cf. \SGARefLink{sga3:i:definition:3-2}{3.2}). A morphism of
\(G\)-\(\mathcal O_S\)-modules is, by definition, a morphism of the
associated \(\mathbf h_G\)-\(\mathbf O_S\)-modules. This gives the
category \((G\text{-}\mathcal O_S\text{-}\mathbf{Mod.})\); we denote by
\((G\text{-}\mathcal O_S\text{-}\mathbf{Mod.q.c.})\) its full
subcategory consisting of those \(G\)-\(\mathcal O_S\)-modules that are
quasi-coherent as \(\mathcal O_S\)-modules.

By \SGARefLink{sga3:i:definition:3-2}{3.2}, giving a \(G\)-\(\mathcal O_S\)-module structure on
\(\mathscr F\) is therefore equivalent to giving a morphism of
\(\widehat{(\mathbf{Sch})}_{/S}\)-groups
\[
  \rho:\mathbf h_G\longrightarrow
    \underline{\operatorname{Aut}}_{\mathbf O_S}
      (\mathbf W(\mathscr F)).
\]

\medskip
\noindent\SGAInlineRefTarget{sga3:i:remark:4-7-1-0}\textbf{Remark 4.7.1.0.}
\footnote{Editor's note: this remark has been added from [DG], II,
\SGARefLink{sga3:ii:section:2}{Section 2}, \SGARefLink{sga3:i:subsection:1-1}{1.1}.}
Since \SGARefLink{sga3:i:proposition:4-6-2}{4.6.2} gives an anti-isomorphism of \(S\)-functors in groups
\[
  (\dagger)\qquad
  \underline{\operatorname{Aut}}_{\mathbf O_S}
    (\mathbf W(\mathscr F))
  \simeq
  \underline{\operatorname{Aut}}_{\mathbf O_S}
    (\mathbf V(\mathscr F)),
\]
it is equivalent to give an \(\mathbf h_G\)-\(\mathbf O_S\)-module
structure on \(\mathbf W(\mathscr F)\) or on \(\mathbf V(\mathscr F)\).
Indeed, let
\[
  \rho:\mathbf h_G\longrightarrow
  \underline{\operatorname{Aut}}_{\mathbf O_S}
    (\mathbf W(\mathscr F))
\]
be as above. For every \(T\to S\) and \(g\in G(T)\), denote by
\(\rho^*(g)\) the image of \(\rho(g)\) under the anti-isomorphism
\((\dagger)\). Then
\[
  \rho^*(gh)=\rho^*(h)\circ\rho^*(g),
\]
so \(\rho^*\) defines a \emph{right}
\(\mathbf h_G\)-\(\mathbf O_S\)-module structure on
\(\mathbf V(\mathscr F)\). Setting
\[
  \rho^{\vee}(g)=\rho^*(g^{-1}),
\]
one obtains an \(\mathbf h_G\)-\(\mathbf O_S\)-module structure on
\(\mathbf V(\mathscr F)\), whose datum is equivalent to that of
\(\rho\).

\medskip
\noindent\SGAInlineRefTarget{sga3:i:remark:4-7-1-1}\textbf{Remark 4.7.1.1.}
One may say that the categories just defined have been constructed by
the following Cartesian squares.

\begin{figure}[H]
  \centering
  \begin{tikzcd}[column sep=3.3em,row sep=3.2em]
  (G\text{-}\mathcal O_S\text{-}\mathbf{Mod.q.c.})
    \arrow[r,hook]
    \arrow[d]
  &
  (G\text{-}\mathcal O_S\text{-}\mathbf{Mod.})
    \arrow[r]
    \arrow[d]
  &
  (\mathbf h_G\text{-}\mathbf O_S\text{-}\mathbf{Mod.})
    \arrow[d,"\mathrm{forget}"]
  \\
  (\mathcal O_S\text{-}\mathbf{Mod.q.c.})
    \arrow[r,hook]
  &
  (\mathcal O_S\text{-}\mathbf{Mod.})
    \arrow[r,"\mathbf W"]
  &
  (\mathbf O_S\text{-}\mathbf{Mod.})
  \qquad .
\end{tikzcd}

  \phantomsection\label{fig:sga3-exp1-p028-rem4711-category-squares}
\end{figure}

The categories \((\mathcal O_S\text{-}\mathbf{Mod.})\) and
\((\mathbf O_S\text{-}\mathbf{Mod.})\) are abelian, but one should
remember that in general the functor \(\mathbf W\) is neither left nor
right exact.

\medskip
\noindent\SGAInlineRefTarget{sga3:i:remark:4-7-1-2}\textbf{Remark 4.7.1.2.}
\footnote{Editor's note: this remark has been added.}
Let \(\mathscr F\) be a \(G\)-\(\mathcal O_S\)-module. The subsheaf of
invariants \(\mathscr F^G\) is defined as follows: for every open subset
\(U\) of \(S\),
\[
  \mathscr F^G(U)=\mathbf W(\mathscr F)^G(U)
   =\bigl\{x\in\mathscr F(U)\mid
       g\cdot x_{S'}=x_{S'}
       \text{ for every }S'\xrightarrow{f}U,
       \ g\in G(S')\bigr\},
\]
where \(x_{S'}\) denotes the image of \(x\) in
\[
  \Gamma(S',f^*\mathscr F)=\Gamma(U,f_*f^*\mathscr F).
\]
One should note that the natural morphism
\[
  \mathbf W(\mathscr F^G)\longrightarrow\mathbf W(\mathscr F)^G
\]
is not an isomorphism in general. For example, let
\(S=\operatorname{Spec}(\mathbb Z)\), and let \(G\) be the constant
group \(\mathbb Z/2\mathbb Z=\{1,\tau\}\) acting on
\(\mathscr F=\mathcal O_S\) by \(\tau\cdot1=-1\). Then
\(\mathscr F^G=0\), whereas, if \(R\) is an \(\mathbb F_2\)-algebra,
\[
  \mathbf W(\mathscr F)^G(\operatorname{Spec}(R))=R.
\]

\medskip
\noindent\SGAInlineRefTarget{sga3:i:paragraph:4-7-2}\textbf{4.7.2.}
From now until the end of \SGARefLink{sga3:i:subsection:4-7}{4.7}, suppose that \(G\) is affine over
\(S\).
By \SGARefLink{sga3:i:proposition:4-6-4}{4.6.4}, giving a morphism of \(S\)-functors
\[
  \rho:\mathbf h_G\longrightarrow
    \underline{\operatorname{End}}_{\mathbf O_S}
      (\mathbf W(\mathscr F))
\]
is equivalent to giving a morphism of \(\mathcal O_S\)-modules
\[
  \mu:\mathscr F\longrightarrow
    \mathscr F\otimes_{\mathcal O_S}\mathscr A(G).
\]
Saying that \(\rho\) is a morphism of
\(\widehat{(\mathbf{Sch})}_{/S}\)-groups is then equivalent to saying
that \(\mu\) satisfies the following axioms.

\medskip
\noindent\textbf{(CM 1)} The following diagram is commutative.

\begin{figure}[H]
  \centering
  \begin{tikzcd}[column sep=5.3em,row sep=3.5em]
  \mathscr F
    \arrow[r,"\mu"]
    \arrow[d,"\mu"']
  &
  \mathscr F\otimes_{\mathcal O_S}\mathscr A(G)
    \arrow[d,"\operatorname{Id}\otimes\Delta"]
  \\
  \mathscr F\otimes_{\mathcal O_S}\mathscr A(G)
    \arrow[r,"\mu\otimes\operatorname{Id}"]
  &
  \mathscr F\otimes_{\mathcal O_S}\mathscr A(G)
    \otimes_{\mathcal O_S}\mathscr A(G)
  \qquad .
\end{tikzcd}

  \phantomsection\label{fig:sga3-exp1-p029-cm1-coaction-square}
\end{figure}

\noindent\textbf{(CM 2)} The following composite is the identity.

\begin{figure}[H]
  \centering
  \begin{tikzcd}[column sep=4.5em]
  \mathscr F
    \arrow[r,"\mu"]
  &
  \mathscr F\otimes\mathscr A(G)
    \arrow[r,"\operatorname{Id}\otimes\varepsilon"]
  &
  \mathscr F\otimes\mathcal O_S
    \arrow[r,"\sim"]
  &
  \mathscr F
  \qquad .
\end{tikzcd}

  \phantomsection\label{fig:sga3-exp1-p029-cm2-counit-composite}
\end{figure}

These axioms (CM 1) and (CM 2) are those of a right \emph{comodule}
structure over the bialgebra \(\mathscr A(G)\).

Put \(\mathscr A=\mathscr A(G)\). If \(\mathscr F\) and
\(\mathscr F'\) are \(\mathscr A\)-comodules, a morphism of comodules
\(f:\mathscr F\to\mathscr F'\) is a morphism of
\(\mathcal O_S\)-modules such that the following diagram is

commutative.

\begin{figure}[H]
  \centering
  \begin{tikzcd}[column sep=4.5em,row sep=3.0em]
  \mathscr F
    \arrow[r,"f"]
    \arrow[d,"\mu_{\mathscr F}"']
  &
  \mathscr F'
    \arrow[d,"\mu_{\mathscr F'}"]
  \\
  \mathscr F\otimes\mathscr A
    \arrow[r,"f\otimes\operatorname{Id}"]
  &
  \mathscr F'\otimes\mathscr A
  \qquad .
\end{tikzcd}

  \phantomsection\label{fig:sga3-exp1-p030-comodule-morphism-square}
\end{figure}

This gives the category \((\mathscr A\text{-}\mathbf{Comod.})\); we
denote by \((\mathscr A\text{-}\mathbf{Comod.q.c.})\) its full
subcategory consisting of those \(\mathscr A\)-comodules that are
quasi-coherent as \(\mathcal O_S\)-modules. We have therefore proved the
following statement.

\medskip
\noindent\SGAInlineRefTarget{sga3:i:proposition:4-7-2}\textbf{Proposition 4.7.2.}
\textit{Let \(G\) be an \(S\)-group affine over \(S\). There are
equivalences of categories
\[
  (G\text{-}\mathcal O_S\text{-}\mathbf{Mod.})
    \simeq(\mathscr A(G)\text{-}\mathbf{Comod.}),
\]
\[
  (G\text{-}\mathcal O_S\text{-}\mathbf{Mod.q.c.})
    \simeq(\mathscr A(G)\text{-}\mathbf{Comod.q.c.}).
\]}

If, in addition, \(S=\operatorname{Spec}(\Lambda)\) is affine and one
sets \(\Lambda[G]=\Gamma(S,\mathscr A(G))\), there is an equivalence of
categories
\[
  (\mathscr A(G)\text{-}\mathbf{Comod.q.c.})
    \simeq(\Lambda[G]\text{-}\mathbf{Comod.}).
\]
\footnote{Editor's note: the preceding sentence has been added.}

Suppose, moreover, that \(\mathscr A=\mathscr A(G)\) is a flat
\(\mathcal O_S\)-module.
\footnote{Editor's note: the following paragraph has been added from
[Se68, \SGARefLink{sga3:i:subsection:1-3}{Section 1.3}].}
Let \(\mathscr E\) be an \(\mathscr A\)-comodule and \(\mathscr F\) an
\(\mathcal O_S\)-submodule of \(\mathscr E\). Since \(\mathscr A\) is
flat over \(\mathcal O_S\), one may identify
\(\mathscr F\otimes\mathscr A\) (respectively,
\(\mathscr F\otimes\mathscr A\otimes\mathscr A\)) with an
\(\mathcal O_S\)-submodule of \(\mathscr E\otimes\mathscr A\)
(respectively, of
\(\mathscr E\otimes\mathscr A\otimes\mathscr A\)). Suppose that
\(\mu_{\mathscr E}\) maps \(\mathscr F\) into
\(\mathscr F\otimes\mathscr A\). Then its restriction
\[
  \mu_{\mathscr F}:\mathscr F\longrightarrow
    \mathscr F\otimes\mathscr A
\]
endows \(\mathscr F\) with an \(\mathscr A\)-comodule structure; one
says that \(\mathscr F\) is a \emph{subcomodule} of \(\mathscr E\).
Passing to the quotient, \(\mu_{\mathscr E}\) defines a morphism of
\(\mathcal O_S\)-modules
\[
  \mathscr E/\mathscr F\longrightarrow
  (\mathscr E/\mathscr F)\otimes\mathscr A,
\]
which endows \(\mathscr E/\mathscr F\) with an
\(\mathscr A\)-comodule structure. If
\(f:\mathscr E\to\mathscr E'\) is a morphism of
\(\mathscr A\)-comodules, then \(\operatorname{Ker}f\) (respectively,
\(\operatorname{Im}f\)) is an \(\mathscr A\)-subcomodule of
\(\mathscr E\) (respectively, of \(\mathscr E'\)), and \(f\) induces
an isomorphism of \(\mathscr A\)-comodules
\[
  \mathscr E/\operatorname{Ker}f
    \xrightarrow{\ \sim\ }
  \operatorname{Im}f.
\]
Moreover, if \(\mathscr E\) and \(\mathscr E'\) are quasi-coherent
\(\mathcal O_S\)-modules, then so are \(\operatorname{Ker}f\) and
\(\operatorname{Im}f\). Consequently,
\((\mathscr A\text{-}\mathbf{Comod.})\) and
\((\mathscr A\text{-}\mathbf{Comod.q.c.})\) are abelian categories.

\medskip
\noindent\SGAInlineRefTarget{sga3:i:corollary:4-7-2-1}\textbf{Corollary 4.7.2.1.}
\textit{Suppose that \(G\) is affine and flat over \(S\). Then the
category \((G\text{-}\mathcal O_S\text{-}\mathbf{Mod.q.c.})\)
(respectively, \((G\text{-}\mathcal O_S\text{-}\mathbf{Mod.})\)),
equivalent to the category of \(\mathscr A(G)\)-comodules that are
quasi-coherent over \(\mathcal O_S\) (respectively, of all
\(\mathscr A(G)\)-comodules), is abelian.}

\medskip
\noindent\SGAInlineRefTarget{sga3:i:paragraph:4-7-3}\textbf{4.7.3.}
Suppose now that \(G\) is a \emph{diagonalizable group}, that is, that
\(\mathscr A(G)\) is the algebra of a commutative group \(M\) over the
sheaf of rings \(\mathcal O_S\). If \(\mathscr F\) is an
\(\mathcal O_S\)-module, then
\[
  \mathscr F\otimes\mathscr A(G)
    =\coprod_{m\in M}\mathscr F\otimes m\mathcal O_S.
\]
Giving a morphism of \(\mathcal O_S\)-modules
\[
  \mu:\mathscr F\longrightarrow\mathscr F\otimes\mathscr A(G)
\]

is therefore equivalent to giving \(\mathcal O_S\)-endomorphisms
\((\mu_m)_{m\in M}\) of \(\mathscr F\) such that, for every section
\(x\) of \(\mathscr F\) on an open subset of \(S\), the family
\((\mu_m(x))\) is a section of the direct sum
\(\coprod_{m\in M}\mathscr F\). This means that on every sufficiently
small open subset, only finitely many restrictions of the
\(\mu_m(x)\) are nonzero.

For \(\mu\) defined by
\[
  \mu(x)=\sum_{m\in M}\mu_m(x)\otimes m,
\]
condition (CM 1), respectively (CM 2), holds if and only if
\[
  \mu_m\circ\mu_{m'}=\delta_{mm'}\mu_m,
  \qquad
  \left(\text{respectively, }\sum_{m\in M}\mu_m
    =\operatorname{Id}_{\mathscr F}\right).
\]
Thus the \(\mu_m\) are pairwise orthogonal projectors whose sum is the
identity. We have proved the following result.

\medskip
\noindent\SGAInlineRefTarget{sga3:i:proposition:4-7-3}\textbf{Proposition 4.7.3.}
\textit{If \(G=D_S(M)\) is a diagonalizable \(S\)-group, the category
of quasi-coherent \(G\)-\(\mathcal O_S\)-modules (respectively, of all
\(G\)-\(\mathcal O_S\)-modules) is equivalent to the category of
quasi-coherent \(\mathcal O_S\)-modules (respectively, of all
\(\mathcal O_S\)-modules) graded by \(M\).}

\medskip
\noindent\textbf{Remark.}
If \(\mathscr F\) has an \(\mathcal O_S\)-algebra structure preserved
by the operations of \(G\), then the grading of \(\mathscr F\) is an
algebra grading. More precisely:

\medskip
\noindent\SGAInlineRefTarget{sga3:i:corollary:4-7-3-1}\textbf{Corollary 4.7.3.1.}
\textit{The functor
\(\mathscr A\mapsto\operatorname{Spec}(\mathscr A)\) induces an
equivalence between the category of quasi-coherent
\(\mathcal O_S\)-algebras graded by \(M\) and the category opposite to
that of \(S\)-schemes affine over \(S\) with \(S\)-operator group
\(G=D_S(M)\).}

\medskip
\noindent\SGAInlineRefTarget{sga3:i:proposition:4-7-4}\textbf{Proposition 4.7.4.}
\textit{Let \(G\) be a diagonalizable \(S\)-group. If
\[
  0\longrightarrow\mathscr F_1\longrightarrow\mathscr F_2
    \longrightarrow\mathscr F_3\longrightarrow0
\]
is an exact sequence of quasi-coherent \(G\)-\(\mathcal O_S\)-modules
that splits as a sequence of \(\mathcal O_S\)-modules, then it also
splits as a sequence of \(G\)-\(\mathcal O_S\)-modules.}

Indeed, if \(G=D_S(M)\), each \(\mathscr F_i\) is graded by the
\((\mathscr F_i)_m\), and for every \(m\in M\) the sequence
\[
  0\longrightarrow(\mathscr F_1)_m\longrightarrow(\mathscr F_2)_m
    \longrightarrow(\mathscr F_3)_m\longrightarrow0
\]
of \(\mathcal O_S\)-modules is split. The preceding proposition then
gives the result.

